\section{Introduction}

This document provides the mathematically rigorous uplift of the A6.x substrate
modules (A–F) and the A7 operator altitude (A7.0–A7.4). Each conceptual module
is translated into formal mathematical objects: symplectic curvature operators,
projection fibrations, unistochastic kernels, sheaf gluing functors, trajectory
invariants, composite algebras, dual-compatibility operators, collapse-surface
repair maps, reprojection functors, and closure theorems.

The result is a unified, altitude-neutral mathematical framework suitable for
verification, peer review, and formal analysis.

\section{Module A: Symplectic Curvature, Torsion, and Drift}

Let $\mathcal{C}$ be the RSVP configuration manifold with symplectic form
$\omega$. Define curvature, torsion, and drift operators:

\[
\mathcal{R} : T\mathcal{C} \to T\mathcal{C}, \qquad
\mathcal{T} : T\mathcal{C} \to T\mathcal{C}, \qquad
\mathcal{D} : T\mathcal{C} \to T\mathcal{C}.
\]

Curvature is defined by:

\[
\mathcal{R}(X,Y) = \nabla_X \nabla_Y - \nabla_Y \nabla_X - \nabla_{[X,Y]}.
\]

Torsion:

\[
\mathcal{T}(X,Y) = \nabla_X Y - \nabla_Y X - [X,Y].
\]

Drift is defined as the Hamiltonian vector field:

\[
\mathcal{D} = X_H = \omega^{-1}(dH).
\]

\section{Module B: Projection Fibrations and Modal Adjacency}

Define the TARTAN projection as a fibration:

\[
\pi : \mathcal{C} \to \mathcal{T},
\]
where $\mathcal{T}$ is the tile manifold.

Modal adjacency is defined by:

\[
p \sim q \iff \exists U \subseteq \mathcal{T} : p,q \in \pi^{-1}(U).
\]

Spectral continuity is defined by:

\[
\mathrm{Spec}(\pi(X)) \text{ continuous in } X.
\]

\section{Module C: Unistochastic Kernels}

Define the amplitude:

\[
U_{ij} = \frac{1}{Z} \exp(i\,\Omega[\gamma_{i\to j}]).
\]

Define the unistochastic kernel:

\[
P_{ij} = |U_{ij}|^2.
\]

\begin{theorem}
If RSVP evolution is symplectic and TARTAN preserves phase modulo coarse
graining, then $P$ is unistochastic.
\end{theorem}

\section{Module D: Sheaf-Cognition Gluing Functors}

Let $\mathcal{O}$ be a sheaf over $M$.

Define gluing functor:

\[
G : \mathrm{Cov}(M) \to \mathrm{Sh}(M)
\]

such that:

\[
G(\{U_i\}) = s \text{ iff } s|_{U_i} = s_i.
\]

Semantic and provenance lifts are defined as natural transformations:

\[
\eta_{\mathrm{sem}} : \mathcal{O} \Rightarrow \mathcal{O}, \qquad
\eta_{\mathrm{prov}} : \mathcal{O} \Rightarrow \mathcal{O}.
\]

\section{Module E: Trajectory Invariants}

Let $\gamma : [0,1] \to \mathcal{C}$ be a smooth trajectory.

Define invariant:

\[
I(\gamma) = \int_0^1 \Phi(\gamma(t))\, dt.
\]

Define stability envelope:

\[
E(\gamma) = \sup_{t} \|\dot{\gamma}(t)\|.
\]

\section{Module F: Composite Substrate Algebra}

Define composite algebra:

\[
\mathcal{A}_{\mathrm{comp}} =
\mathcal{A}_{\mathrm{RSVP}} \oplus
\mathcal{A}_{\mathrm{TARTAN}} \oplus
\mathcal{A}_{\mathrm{uni}} \oplus
\mathcal{A}_{\mathrm{sheaf}} \oplus
\mathcal{A}_{\mathrm{traj}}.
\]

Define composite operator:

\[
\mathcal{O}(p) =
\mathcal{R}(p) \oplus
\mathcal{D}(p) \oplus
\mathcal{S}(p) \oplus
\mathcal{X}(p) \oplus
\mathcal{E}(p).
\]

\section{A7.0: Composite Operator Algebra}

Define Lie bracket:

\[
[\mathcal{O}_1,\mathcal{O}_2] =
[\mathcal{R}_1,\mathcal{R}_2] \oplus
[\mathcal{D}_1,\mathcal{D}_2] \oplus
[\mathcal{S}_1,\mathcal{S}_2] \oplus
[\mathcal{X}_1,\mathcal{X}_2] \oplus
[\mathcal{E}_1,\mathcal{E}_2].
\]

\section{A7.1: Dual-Compatibility Operators}

Define dual substrate:

\[
\mathcal{A}^\star_{\mathrm{comp}} = \mathrm{Hom}(\mathcal{A}_{\mathrm{comp}},\mathbb{R}).
\]

Define dual-compatibility operator:

\[
\mathcal{T}(p)(\mathcal{O}) = \langle \mathcal{O}, K(p) \rangle.
\]

\section{A7.2: Collapse-Surface Repair Maps}

Define collapse surface $\Sigma$.

Define repair map:

\[
\mathcal{R} : \Sigma \to \Sigma
\]

such that:

\[
\mathcal{R}(p) = p.
\]

\section{A7.3: Reprojection Functors}

Define reprojection functor:

\[
\mathcal{P} : \mathcal{A}_{\mathrm{comp}} \to \mathcal{A}_{\mathrm{comp}}
\]

such that:

\[
\mathcal{P}(\mathcal{O}) = \mathcal{O}.
\]

\section{A7.4: Completion Conditions}

Define closure conditions:

\[
\mathcal{T}(p)(\mathcal{O}(p)) = \mathcal{O}(p),
\]
\[
\ker(\mathcal{O}(p)) = \ker(\mathcal{T}(p)(\mathcal{O}(p))),
\]
\[
\mathrm{Spec}(\mathcal{O}(p)) = \mathrm{Spec}^\star(\mathcal{T}(p)).
\]

\section{Conclusion}

This document provides the full mathematical translation of the A6.x substrate
modules and the A7 operator altitude into a unified, academically defensible
framework.
